\problemname{Eating Chocolate}

\noindent
Olle and Lisa have received $N$ chocolate bars of various shapes from their parents, which they now plan to eat together (the chocolate bars, not the parents).
Each chocolate bar is a rectangle consisting of $R_i \times C_i$ squares.
The chocolate bars are placed on a table, and to eat them, Olle and Lisa will take turns making the following moves:
\begin{itemize}
  \item Choose a chocolate bar from the table.
  \item Break it into two parts, so that both parts are rectangles with integer side lengths.
  \item Eat the smaller part.
  \item Place the remaining part back on the table.
\end{itemize}
If the chosen chocolate bar has a size of $1 \times 1$, they eat it directly and do not place anything back on the table.
They have already decided that Lisa will start.


Due to being siblings, they will both want to eat as much chocolate as possible.
In each move, they will choose the chocolate bar that, after they have split it and put the bigger part back on the table,
results in them eating as many squares as possible in that move.
If multiple such choices are possible, they will arbitrarily choose one of them such that their future chocolate consumtion is maximized.
Additionally, Olle cares a lot about fairness and would like the total number of squares eaten by both to be as equal as possible.
To achieve this goal, before they start he can choose to discard at most one chocolate bar into the trash can without Lisa noticing.

Let $V_O$ and $V_L$ be the number of squares eaten by Olle and Lisa, respectively, given that they both eat in a greedy and logical manner.  
Your task is to determine the minimal difference, that is, $|V_O - V_L|$,  
after Olle has discarded at most one chocolate bar.

\section*{Input}
The first line contains the number $N$ ($2 \leq N \leq 10^5$), the number of chocolate bars on the table.

This is followed by $N$ lines, where the $i$:th line contains the numbers $R_i$ and $C_i$ ($1 \leq R_i, C_i \leq 10^6$).

\section*{Output}
Print a single integer, the minimal value of $|V_O - V_L|$ as described above.

\section*{Scoring}
Your solution will be tested on a set of test groups, each worth a number of points. Each test group contains
a set of test cases. To get the points for a test group you need to solve all test cases in the test group.


\noindent
\begin{tabular}{| l | l | p{12cm} |}
  \hline
  \textbf{Group} & \textbf{Points} & \textbf{Constraints} \\ \hline
  $1$    & $10$       & All chocolate bars have the same size. \\ \hline
  $2$    & $16$       & $R_i, C_i \leq 10$ \\ \hline
  $3$    & $15$       & $N \leq 50$ \\ \hline
  $4$    & $19$       & $N \leq 500$ \\ \hline
  $5$    & $40$       & No additional constraints. \\ \hline
\end{tabular}

\section*{Explanation of Sample 1}
Olle discards the chocolate bar of size $2 \times 2$.  
After that, they will each eat one square of chocolate twice, so $V_O = V_L$.

\section*{Explanation of Sample 2}
In this test case, the optimal strategy for Olle is to not discard any chocolate bar.  
They will then eat exactly the same amount of chocolate.
